% ==============================================================================
%  3.2  Inferencias sobre la respuesta media
%  3.3  Inferencias sobre una nueva observación
% ==============================================================================
\section{Inferencias sobre la respuesta media}
\label{sec:t3-respuesta-media}

Un objetivo habitual de los modelos de regresión es estimar la respuesta media para un valor
dado $x_h$ de $X$. En el \cref{ej:arboles}, interesa estimar el volumen medio de madera de los
árboles con un determinado diámetro, para decidir si merece la pena talarlos.

Denotamos $\E(Y_h)=\E(Y\cond X=x_h)=\beta_0+\beta_1x_h$. Su \textbf{estimador puntual} es el
valor de la recta ajustada en $x_h$:
\[
  \est{Y}_h = \est{\beta}_0 + \est{\beta}_1x_h .
\]

\begin{teorema}[Distribución de $\est{Y}_h$][teo:t3-dist-Yh]
  \begin{gather*}
    \est{Y}_h \sim \Normal\!\left(\E(Y_h),\ \sigma^2\left(\frac1n+\frac{(x_h-\media{X})^2}{\SSX}\right)\right),
    \\
    \frac{\est{Y}_h-\E(Y_h)}{\raizMSE{\frac1n+\frac{(x_h-\media{X})^2}{\SSX}}} \sim \tdist{n-2}.
  \end{gather*}
\end{teorema}

\begin{demostracion}
  \emph{Normalidad.} $\est{Y}_h$ es combinación lineal de $\est{\beta}_0$ y $\est{\beta}_1$, y por
  tanto de las $Y_i$, que son normales independientes; luego es normal.

  \emph{Insesgadez.}
  \begin{align*}
    \E(\est{Y}_h) &= \E(\est{\beta}_0)+\E(\est{\beta}_1)\,x_h \\
    &= \beta_0+\beta_1x_h \\
    &= \E(Y_h) .
  \end{align*}

  \emph{Varianza.} Usando $\est{\beta}_0=\media{Y}-\est{\beta}_1\media{X}$,
  \begin{align*}
    \est{Y}_h &= \est{\beta}_0+\est{\beta}_1x_h \\
    &= \media{Y}-\est{\beta}_1\media{X}+\est{\beta}_1x_h \\
    &= \media{Y}+\est{\beta}_1(x_h-\media{X}) ,
  \end{align*}
  y, como $x_h$ y $\media{X}$ son constantes,
  \[
    \Var(\est{Y}_h) = \Var(\media{Y}) + (x_h-\media{X})^2\,\Var(\est{\beta}_1)
    + 2(x_h-\media{X})\Cov(\media{Y},\est{\beta}_1) .
  \]
  La covarianza es nula: escribiendo $\media{Y}$ y $\est{\beta}_1$ como combinaciones lineales de
  las $Y_i$, que son independientes con varianza $\sigma^2$, y usando
  la propiedad (i) de los $k_i$ (\cref{prop:t3-ki}),
  \begin{align*}
    \Cov(\media{Y},\est{\beta}_1) &= \Cov\left(\sumi \frac1n\,Y_i,\ \sumi k_iY_i\right) \\
    &= \sumi \frac1n\,k_i\,\Var(Y_i) \\
    &= \frac{\sigma^2}{n}\sumi k_i = 0 .
  \end{align*}
  Por tanto, con $\Var(\media{Y})=\sigma^2/n$ y $\Var(\est{\beta}_1)=\sigma^2/\SSX$
  (\cref{teo:t3-dist-b1}),
  \begin{align*}
    \Var(\est{Y}_h) &= \frac{\sigma^2}{n} + (x_h-\media{X})^2\,\frac{\sigma^2}{\SSX} \\
    &= \sigma^2\left(\frac1n+\frac{(x_h-\media{X})^2}{\SSX}\right).
  \end{align*}
  Al sustituir $\sigma^2$ por $MSE$ se obtiene la distribución $\tdist{n-2}$, por el mismo
  argumento que en el \cref{teo:t3-dist-b1} ($\est{Y}_h$ es función de $\est{\beta}_0$ y
  $\est{\beta}_1$, independientes de $SSE$).
\end{demostracion}

\begin{resumen}[IC de nivel $1-\alpha$ para la respuesta media $\E(Y_h)$]
  \[
    \E(Y_h) \in \left(\est{Y}_h \pm \tq{n-2}{1-\alpha/2}
    \raizMSE{\frac1n+\frac{(x_h-\media{X})^2}{\SSX}}\right).
  \]
\end{resumen}

La varianza de $\est{Y}_h$ es tanto menor cuanto más próximo está $x_h$ a la media muestral
$\media{X}$: los intervalos más precisos se obtienen en el punto $(\media{X},\media{Y})$ y se
ensanchan al alejarse de él.

% ------------------------------------------------------------------------------
\section{Inferencias sobre una nueva observación}
\label{sec:t3-prediccion}

Se quiere ahora predecir el valor de una \textbf{nueva observación} de $Y$, $\Ynew$, cuando
$X=x_h$. Por ejemplo, se mide el diámetro de un árbol concreto, 10.2 pulgadas, y se quiere
predecir el volumen de madera de \emph{ese} árbol.

\begin{itemize}
  \item La nueva observación corresponde a un nuevo individuo, al que se asume aplicable el
    modelo de regresión ajustado.
  \item Hay que distinguir entre la predicción de la respuesta media, $\E(Y_h)$, y la de una
    nueva observación, $\Ynew$:
    \begin{itemize}
      \item $\E(Y_h)$ es la media de la distribución de $Y$ condicionada a $X=x_h$;
      \item $\Ynew$ es un valor de una variable aleatoria con la distribución de $Y$
        condicionada a $X=x_h$, y difícilmente coincidirá con $\E(Y_h)$.
    \end{itemize}
\end{itemize}
El \textbf{estimador puntual} es el mismo que para la respuesta media,
\[
  \est{Y}_{h(\mathrm{new})} = \est{Y}_h = \est{\beta}_0 + \est{\beta}_1x_h ,
\]
pero el error que se comete es distinto, y por tanto también su distribución.

\paragraph{Parámetros conocidos.} Si $\beta_0$, $\beta_1$ y $\sigma^2$ fueran conocidos, la única
fuente de incertidumbre sería la variabilidad de $Y$ alrededor de su media:
$\Ynew\sim\Normal(\beta_0+\beta_1x_h,\sigma^2)$, y un intervalo de nivel $1-\alpha$ sería
$\E(Y_h)\pm z_{1-\alpha/2}\,\sigma$.

\begin{ejemplo}[Volumen de madera con parámetros conocidos]
  Supongamos $\beta_0=-45.57$, $\beta_1=6.916$ y $\sigma^2=36.587$ ($\sigma=6.049$).
  Para un árbol
  de diámetro $x_h=10.2$ pulgadas, el volumen medio es
  $\E(Y_h)=-45.57+6.916\cdot10.2=24.97$ pies cúbicos. Si el modelo es adecuado (errores normales,
  independientes y con varianza constante), el volumen de ese árbol sigue una
  $\Normal(24.97,\ 36.587)$, y un intervalo del 95\,\% para él es
  \[
    24.97 \pm 1.96\cdot6.049 = (13.11,\ 36.83).
  \]
\end{ejemplo}

\paragraph{Parámetros desconocidos.} En la práctica los parámetros se estiman, y al estimar
$\E(Y_h)$ se comete un error: el centro del intervalo para $\Ynew$ es a su vez incierto (no se
conoce con certeza la localización de la distribución de $Y$). Hay dos fuentes de
variabilidad:
\begin{enumerate}
  \item la asociada a la estimación de la respuesta media $\E(Y_h)$;
  \item la asociada a la distribución de $Y$ alrededor de su media.
\end{enumerate}

\begin{figure}[H]
  \centering
  \begin{tikzpicture}[x=1cm, y=2.2cm]
    \draw[-{Stealth}] (-4.8,0) -- (4.8,0);
    \foreach \c in {-2,2} {
      \draw[thick, domain=\c-2.6:\c+2.6, samples=60] plot (\x, {exp(-(\x-\c)^2/1.3)});
      \draw[dashed] (\c,0) -- (\c,1);
      \draw (\c-1.5,0) -- (\c-1.5,1.35);
      \draw (\c+1.5,0) -- (\c+1.5,1.35);
      \draw[{Stealth}-{Stealth}] (\c-1.5,1.25) -- (\c+1.5,1.25);
      \node[font=\scriptsize, fill=white, align=center] at (\c,1.25)
        {límites de predicción\\ si $\E(Y_h)$ está aquí};
    }
    \draw (0,0.05) -- (0,-0.05) node[below, font=\small] {$\est{Y}_h$};
    \draw[{Stealth}-{Stealth}] (-2,-0.45) -- (2,-0.45)
      node[midway, below, font=\scriptsize] {límites de confianza para $\E(Y_h)$};
  \end{tikzpicture}
  \caption{La predicción de una nueva observación combina la incertidumbre sobre
    $\E(Y_h)$ con la variabilidad de $Y$ alrededor de su media.}
\end{figure}

La cantidad relevante es el \textbf{error de predicción} $\Ynew-\est{Y}_h$.
Como $\Ynew$ es
independiente de las observaciones de la muestra, y por tanto de $\est{Y}_h$,
\begin{align*}
  \Var(\Ynew-\est{Y}_h) &= \Var(\Ynew) + \Var(\est{Y}_h)
  = \sigma^2 + \sigma^2\left(\frac1n+\frac{(x_h-\media{X})^2}{\SSX}\right) \\
  &= \sigma^2\left(1+\frac1n+\frac{(x_h-\media{X})^2}{\SSX}\right),
\end{align*}
y el error de predicción tiene media $\E(Y_h)-\E(Y_h)=0$. Es normal, por ser diferencia de
normales independientes, y es independiente de $SSE$, ya que $\Ynew$ es independiente de la
muestra y $\est{Y}_h$ lo es de $SSE$; por ello, al sustituir $\sigma^2$ por $MSE$ se obtiene de
nuevo una $\tdist{n-2}$.

\begin{teorema}[Intervalo de predicción][teo:t3-prediccion]
  \begin{gather*}
    \Ynew-\est{Y}_h \sim \Normal\!\left(0,\ \sigma^2\left(1+\frac1n+\frac{(x_h-\media{X})^2}{\SSX}\right)\right),
    \\
    \frac{\Ynew-\est{Y}_h}{\raizMSE{1+\frac1n+\frac{(x_h-\media{X})^2}{\SSX}}} \sim \tdist{n-2},
  \end{gather*}
  y el \textbf{intervalo de predicción} de nivel $1-\alpha$ para $\Ynew$ es
  \[
    \Ynew \in \left(\est{Y}_h \pm \tq{n-2}{1-\alpha/2}
    \raizMSE{1+\frac1n+\frac{(x_h-\media{X})^2}{\SSX}}\right).
  \]
\end{teorema}

El intervalo para una nueva observación tiene el mismo centro que el de la respuesta media, pero
mayor varianza. Por tanto, el IC para la respuesta media está siempre contenido en el intervalo de
predicción.

\subsection{Media de \texorpdfstring{$m$}{m} nuevas observaciones}

En ocasiones interesa predecir la media $\Ymnew$ de $m$ nuevas observaciones de $Y$ para
$X=x_h$. Se asume que el modelo es aplicable a los $m$ nuevos individuos, que son por tanto
independientes. Las fuentes de variabilidad son la estimación de $\E(Y_h)$ y la variabilidad de
la media de $m$ observaciones de $Y\cond X=x_h$, con $\Var(\Ymnew)=\sigma^2/m$. Así,
\begin{align*}
  \Var(\Ymnew-\est{Y}_h) &= \frac{\sigma^2}{m} + \sigma^2\left(\frac1n+\frac{(x_h-\media{X})^2}{\SSX}\right) \\
  &= \sigma^2\left(\frac1m+\frac1n+\frac{(x_h-\media{X})^2}{\SSX}\right).
\end{align*}

\begin{resumen}[Intervalo de predicción para la media de $m$ nuevas observaciones]
  \begin{gather*}
    \frac{\Ymnew-\est{Y}_h}{\raizMSE{\frac1m+\frac1n+\frac{(x_h-\media{X})^2}{\SSX}}}\sim\tdist{n-2},
    \\
    \Ymnew \in \left(\est{Y}_h \pm \tq{n-2}{1-\alpha/2}
    \raizMSE{\frac1m+\frac1n+\frac{(x_h-\media{X})^2}{\SSX}}\right).
  \end{gather*}
  Con $m=1$ se recupera el intervalo de predicción del \cref{teo:t3-prediccion}.
\end{resumen}

\subsection{Predicción inversa}

En ocasiones interesa \textbf{predecir el valor de $X$} a partir de un valor observado de $Y$,
utilizando el modelo de $Y$ sobre $X$.

\begin{ejemplo}[Calibración de un aparato de medida][ej:t3-calibracion]
  Se realiza un experimento con 12 individuos para evaluar un aparato que mide la concentración
  de azúcar (galactosa) en sangre. A partir de concentraciones conocidas ($X$) se registran las
  mediciones del aparato ($Y$) y se ajusta el modelo
  \[
    \est{Y} = -0.1 + 1.017\,X .
  \]
  Para un nuevo individuo solo se dispone de la medición del aparato, $\Ynew$, y se quiere
  estimar su verdadera concentración $\Xnew$.
\end{ejemplo}

Despejando en la recta ajustada se obtiene
$\est{X}=(Y+0.1)/1.017=0.098+0.983\,Y$. Esta recta \emph{no} es la recta de regresión de $X$
sobre $Y$: la de $X$ sobre $Y$ se obtendría minimizando las desviaciones en la dirección de $X$,
con pendiente $SS_{XY}/SS_Y$, mientras que al despejar se obtiene pendiente
$1/\est{\beta}_1=\SSX/SS_{XY}$; ambas coinciden solo si $r_{XY}^2=1$. Además, en este problema $X$
es la variable controlada (no aleatoria) e $Y$ la aleatoria, por lo que el modelo adecuado es el
de $Y$ sobre $X$. La estimación de $X$ se hace despejando en él, lo que se denomina
\textbf{predicción inversa}.

\begin{definicion}[Predicción inversa][def:t3-prediccion-inversa]
  Conocido el modelo estimado $\est{Y}=\est{\beta}_0+\est{\beta}_1X$ y observado $\Ynew$ en un
  nuevo individuo, el estimador puntual de $\Xnew$ es
  \[
    \est{X}_{h(\mathrm{new})} = \frac{\Ynew-\est{\beta}_0}{\est{\beta}_1}, \qquad \est{\beta}_1\neq0 .
  \]
\end{definicion}

Para hacer inferencia sobre $\Xnew$ se usan las siguientes propiedades (aproximadas, ya que
$\est{X}_{h(\mathrm{new})}$ es un cociente de variables aleatorias):
\begin{itemize}
  \item es aproximadamente normal;
  \item es aproximadamente insesgado:
    \[
      \E\big(\est{X}_{h(\mathrm{new})}\big)\approx
      \frac{\E(\Ynew)-\E(\est{\beta}_0)}{\E(\est{\beta}_1)}
      = \frac{\beta_0+\beta_1\Xnew-\beta_0}{\beta_1} = \Xnew ;
    \]
  \item su varianza es aproximadamente
    $\Var\big(\est{X}_{h(\mathrm{new})}\big)\approx
    \dfrac{\sigma^2}{\beta_1^2}\left(1+\dfrac1n+\dfrac{(\Xnew-\media{X})^2}{\SSX}\right)$.
\end{itemize}

\begin{resumen}[IC aproximado de nivel $1-\alpha$ para $\Xnew$]
  \[
    \Xnew \in \left(\est{X}_{h(\mathrm{new})} \pm \tq{n-2}{1-\alpha/2}
    \sqrt{\frac{MSE}{\est{\beta}_1^2}\left(1+\frac1n+
    \frac{\big(\est{X}_{h(\mathrm{new})}-\media{X}\big)^2}{\SSX}\right)}\,\right).
  \]
\end{resumen}
